\begin{textsample}{POS Dim 4 – human – Score 5.00 – rj/rj\_ef\_linguagens\_lp\_1\_p2.txt}  \label{ex:f4_pos_011}
A presença de um trabalho pautado na diversidade \textbf{textual} requer um lugar de destaque no cotidiano escolar, pois através dele os estudantes iniciam o processo de autonomia, pois a \textbf{leitura} e a \textbf{escrita} são, incontestavelmente, o caminho para que estejam inseridos em nossa sociedade.
Durante o trabalho pedagógico com a Língua Portuguesa o enfoque do professor \textbf{precisa} ser o \textbf{texto} em seus variados contextos, o que também significa inserir as variantes locais em seu dia a dia escolar, de maneira que o \textbf{texto}, \textbf{oral}, escrito ou multimodal, passe a ser o centro das atividades a serem desenvolvidas no espaço escolar, permitindo que haja uma troca permanente de ideias, questionamentos e posicionamentos, respeitando assim as individualidades e necessidades dos envolvidos no processo.

% matched lemmas: escrita, leitura, oral, precisar, texto, textual
\end{textsample}
